% Options for packages loaded elsewhere
\PassOptionsToPackage{unicode}{hyperref}
\PassOptionsToPackage{hyphens}{url}
\documentclass[
  ignorenonframetext,
]{beamer}
\newif\ifbibliography
\usepackage{pgfpages}
\setbeamertemplate{caption}[numbered]
\setbeamertemplate{caption label separator}{: }
\setbeamercolor{caption name}{fg=normal text.fg}
\beamertemplatenavigationsymbolsempty
% remove section numbering
\setbeamertemplate{part page}{
  \centering
  \begin{beamercolorbox}[sep=16pt,center]{part title}
    \usebeamerfont{part title}\insertpart\par
  \end{beamercolorbox}
}
\setbeamertemplate{section page}{
  \centering
  \begin{beamercolorbox}[sep=12pt,center]{section title}
    \usebeamerfont{section title}\insertsection\par
  \end{beamercolorbox}
}
\setbeamertemplate{subsection page}{
  \centering
  \begin{beamercolorbox}[sep=8pt,center]{subsection title}
    \usebeamerfont{subsection title}\insertsubsection\par
  \end{beamercolorbox}
}
% Prevent slide breaks in the middle of a paragraph
\widowpenalties 1 10000
\raggedbottom
\AtBeginPart{
  \frame{\partpage}
}
\AtBeginSection{
  \ifbibliography
  \else
    \frame{\sectionpage}
  \fi
}
\AtBeginSubsection{
  \frame{\subsectionpage}
}
\usepackage{iftex}
\ifPDFTeX
  \usepackage[T1]{fontenc}
  \usepackage[utf8]{inputenc}
  \usepackage{textcomp} % provide euro and other symbols
\else % if luatex or xetex
  \usepackage{unicode-math} % this also loads fontspec
  \defaultfontfeatures{Scale=MatchLowercase}
  \defaultfontfeatures[\rmfamily]{Ligatures=TeX,Scale=1}
\fi
\usepackage{lmodern}
\ifPDFTeX\else
  % xetex/luatex font selection
\fi
% Use upquote if available, for straight quotes in verbatim environments
\IfFileExists{upquote.sty}{\usepackage{upquote}}{}
\IfFileExists{microtype.sty}{% use microtype if available
  \usepackage[]{microtype}
  \UseMicrotypeSet[protrusion]{basicmath} % disable protrusion for tt fonts
}{}
\makeatletter
\@ifundefined{KOMAClassName}{% if non-KOMA class
  \IfFileExists{parskip.sty}{%
    \usepackage{parskip}
  }{% else
    \setlength{\parindent}{0pt}
    \setlength{\parskip}{6pt plus 2pt minus 1pt}}
}{% if KOMA class
  \KOMAoptions{parskip=half}}
\makeatother
\setlength{\emergencystretch}{3em} % prevent overfull lines
\providecommand{\tightlist}{%
  \setlength{\itemsep}{0pt}\setlength{\parskip}{0pt}}
% Slide layout defaults for warning-clean scientific decks.
\usepackage{etoolbox}
\IfFileExists{xurl.sty}{\usepackage{xurl}}{}
\IfFileExists{seqsplit.sty}{\usepackage{seqsplit}}{\newcommand{\seqsplit}[1]{#1}}
\protected\def\breakseq#1{\seqsplit{#1}}
\protected\def\breaktt#1{\begingroup\ttfamily\seqsplit{#1}\endgroup}
\setlength{\emergencystretch}{6em}
\tolerance=5000
\hbadness=10000
\hfuzz=1pt
\setlength{\tabcolsep}{2pt}
\AtBeginEnvironment{longtable}{\tiny\renewcommand{\arraystretch}{0.86}\setlength{\tabcolsep}{1pt}}
\AtBeginEnvironment{tabular}{\tiny\renewcommand{\arraystretch}{0.86}\setlength{\tabcolsep}{1pt}}

% Natbib and cross-reference fallbacks — slides don't load natbib
% or cleveref, but combined-PDF manuscript prose may emit these
% commands. The fallback renders citations as a bracketed key list
% and cross-references as detokenized labels so slides don't fail on
% undefined control sequences or raw underscores. \providecommand is
% a no-op if packages load later.
\providecommand{\citep}[1]{[#1]}
\providecommand{\citet}[1]{#1}
\providecommand{\citealp}[1]{#1}
\providecommand{\citeauthor}[1]{#1}
\providecommand{\citeyear}[1]{#1}
\providecommand{\cref}[1]{\texttt{\detokenize{#1}}}
\providecommand{\Cref}[1]{\texttt{\detokenize{#1}}}

\newtheorem{proposition}{Proposition}
\newtheorem{hypothesis}{Hypothesis}
\usepackage{bookmark}
\IfFileExists{xurl.sty}{\usepackage{xurl}}{} % add URL line breaks if available
\urlstyle{same}
% fallback for those not using the hyperref driver hyperxmp:
\makeatletter
\@ifundefined{xmpquote}{\newcommand{\xmpquote}[1]{#1}}{}
\makeatother
\hypersetup{
  hidelinks,
  pdfcreator={LaTeX via pandoc}}

\author{\texorpdfstring{}{}}
\date{}

\begin{document}

\section{Abstract}\label{abstract}

\begin{frame}[allowframebreaks]{Abstract}
No prior automated system tracks hypothesis-level evidence mapping
across the full Active Inference and Free Energy Principle (FEP)
literature. Manual synthesis cannot keep pace with a field that has
grown at a compound annual rate of 20.36\% across 2000--2026 (empirical
publication span 2005--2026), and the FEP's theoretical generality has
invited falsifiability critiques that only hypothesis-specific triage
can address. Building on pioneering systematic manual annotation paired
with ontology-based analysis at the scale of hundreds of papers, we
present a computational meta-analysis framework that automates and
scales this approach. The pipeline retrieves literature from arXiv,
Semantic Scholar, and OpenAlex, deduplicating \(N = 817\) papers
(inclusion from 2000 onward) via a canonical identifier hierarchy (DOI
\(>\) arXiv ID \(>\) Semantic Scholar ID \(>\) OpenAlex ID). It
classifies papers into a three-tier taxonomy spanning eight categories:
A (Core Theory), B (Tools \& Translation), and C (Application Domains).
An LLM-powered extraction system evaluates each abstract against eight
core hypotheses, separating (1) explicit source claims, (2) whether the
abstract supplies evidence, and (3) pipeline hypothesis triage---each
encoded in structured nanopublications with full extraction provenance.

All extracted assertions are machine-generated; hypothesis scores report
citation-weighted evidence mapping and triage, not scientific
confirmation. In the absence of a human-annotated gold set, a stratified
agreement study (\(n = 200\)) compares pipeline triage against an
independent, deterministic keyword-rule reference: pipeline precision
0.000 and recall 0.000 against the rule reference, with
reference--pipeline direction agreement of only \(\kappa = -0.055\).
This low agreement is the point---absolute triage labels are unreliable,
so we report relative rankings and trajectories rather than point
estimates; human gold-standard validation remains future work.

The resulting evidence landscape reveals a field where application
domains (Domain C, 64.0\%) collectively dominate the corpus, with tools
development (Domain B, 20.8\%) and core theory (Domain A, 15.2\%)
rounding out the taxonomy. Non-negative matrix factorization identifies
8 latent topics that cross-cut the keyword taxonomy, and citation
network analysis exposes a sparse yet structured graph (2,176
intra-corpus edges out of 29,323 total outgoing
references---\textbf{only 7.4\% reference resolution}) anchored by
pronounced hub papers. Hypothesis triage scores cluster into tiers for
relative comparison; absolute magnitudes are inflated by publication
bias and linguistic asymmetry, so rankings and temporal trajectories are
more reliable than point estimates. This work provides a reusable
architecture for \emph{living literature reviews}---continuously updated
knowledge graphs that map citation-weighted evidence and support
hypothesis triage across rapidly evolving fields.

\textbf{Keywords:} Active Inference, Free Energy Principle,
meta-analysis, knowledge graph, nanopublications, bibliometrics,
hypothesis scoring, LLM extraction, computational neuroscience
\end{frame}

\end{document}
